%!TEX root = ../thesis.tex

\chapter{Model Architecture and Training Objective}\label{app:ez-details}

% TODO: The parameterization that \cref{chap:efficientzero} leaves out on purpose,
% at the level of detail needed to reimplement it:
%   - the bodies of the five networks of \cref{sec:ez-architecture}: block types,
%     depths and widths, and the recurrent value-prefix head,
%   - the projection $P_1$ and prediction $P_2$ heads that \cref{eq:ez-consistency}
%     refers to,
%   - the categorical representation of value and value prefix: support size and
%     the invertible scaling transform
%     \todo{cite: value-scaling transform (Pohlen et al. 2018) --- not in references/*.bib yet},
%   - the weighted total loss behind \cref{sec:ez-loss}. The reference
%     configuration weights policy 1.0, value 0.25, value prefix 1.0 and
%     consistency 2.0 (\code{efficientzero/src/learner/mod.rs}); the value term is
%     down-weighted because its reanalyzed bootstrap targets are the noisiest signal,
%   - gradient scaling across the unroll steps,
%   - how the value loss becomes the replay priority \cite{schaul_prioritized_2016}.
%
% Boundary with \cref{app:configuration}: the formulas belong here, the frozen
% numeric values belong in the hyperparameter table there.

\chapter{Frozen Configuration}\label{app:configuration}

\Cref{tab:app-ez-params} lists the EfficientZero parameters every run of this thesis uses.

\begin{table}[htbp]
  \centering
  \caption[EfficientZero parameters]{%
    The EfficientZero parameters of this work next to the Atari configuration of the reference
    implementation~\cite{ye_mastering_2021}.
    The three search parameters $\xi$, $\rho$ and $T$ are the exploration change of
    \cref{sec:methodology-setup}; every other difference follows from the scale of the targets.
  }
  \label{tab:app-ez-params}
  \footnotesize
  \begin{tabularx}{\textwidth}{@{}ll>{\raggedright\arraybackslash}X>{\raggedright\arraybackslash}X@{}}
    \toprule
    Parameter & & This work & EfficientZero (Atari) \\
    \midrule
    \multicolumn{4}{@{}l}{\emph{Search}} \\
    Simulations per search
      & & 100 & 50 \\
    pUCT constants
      & $c_1$, $c_2$ & 1.25, \num{19652} & 1.25, \num{19652} \\
    Discount
      & $\gamma$ & 0.99 & 0.988 \\
    Value-prefix LSTM horizon
      & & 5 & 5 \\
    Root Dirichlet concentration
      & $\xi$ & 1.0 & 0.3 \\
    Root noise fraction
      & $\rho$ & 0.4 & 0.25 \\
    Visit-count temperature
      & $T$ & $1 \to 0.5 \to 0.25$ at \SI{20}{\percent} and \SI{40}{\percent} of the run
      & $1 \to 0.5 \to 0.25$ at \SI{50}{\percent} and \SI{75}{\percent} of training \\
    \addlinespace
    \multicolumn{4}{@{}l}{\emph{Targets and replay}} \\
    Unroll steps
      & $K$ & 5 & 5 \\
    TD steps
      & $n$ & 5 & 5 \\
    Off-policy correction of $n$
      & & scaled by the buffer capacity & $0.3 \times$ training steps \\
    Reanalyzed policy targets
      & & \SI{99}{\percent} of a batch & \SI{99}{\percent} of a batch \\
    Replay capacity
      & & \num{10000} trajectories & \num{1000000} transitions \\
    Priority and importance-sampling exponents
      & & 0.6, 0.4 & 0.6, 0.4 \\
    \addlinespace
    \multicolumn{4}{@{}l}{\emph{Loss weights}} \\
    Policy, value, value-prefix, consistency
      & $\lambda_p$, $\lambda_v$, $\lambda_x$, $\lambda_c$ & 1, 0.25, 1, 2 & 1, 0.25, 1, 2 \\
    \addlinespace
    \multicolumn{4}{@{}l}{\emph{Optimization}} \\
    Optimizer
      & & AdamW & SGD, momentum 0.9 \\
    Learning rate
      & & \num{1e-3}, constant & 0.2, warm-up and step decay \\
    Weight decay
      & & \num{1e-4} & \num{1e-4} \\
    Gradient-norm clip
      & & 5 & 5 \\
    Batch size
      & & 128 & 256 \\
    \addlinespace
    \multicolumn{4}{@{}l}{\emph{Model}} \\
    Trunk
      & & residual MLP, latent width 128 & convolutional ResNet, 64 channels \\
    Value and value-prefix support
      & & 71 bins on $[-30, 30]$, resized per target & 601 bins on $[-300, 300]$ \\
    \bottomrule
  \end{tabularx}
\end{table}

% TODO: \cref{tab:app-ez-params} covers the EfficientZero parameters only. What is
% still missing is the complete hyperparameter table referenced from
% \cref{sec:methodology-hyperparameters}, annotated by justification tier (inherited,
% derived, sensitivity-analyzed) so the main text can cite it instead of reproducing
% it. Include the commit hash and the date the configuration was frozen.

\chapter{Action and Observation Encodings}\label{app:encodings}

% TODO: The full encodings that \cref{sec:input-action-space,sec:corpus-action-space}
% and the observation sections summarize: the packing arithmetic for the composite
% corpus action, the per-slot feature layout, and the exact vector widths per target.
% Enough detail to reimplement, kept out of the main text so the design argument
% stays readable.

\chapter{Additional Results}\label{app:results}

% TODO: Per-run curves, the peripheral ablations labelled indicative in
% \cref{sec:results-ablations}, and the diagnostic series that support
% \cref{chap:failure-modes} without belonging in it.
